% =====================================================================
%  KDS-EL 论文完整格式规范
%  参考：EMSWOA (Applied Soft Computing 2025) + ASC/ESWA 惯例
%  目标期刊：Expert Systems with Applications / Applied Soft Computing
% =====================================================================


% ─────────────────────────────────────────────────────────────────────
% 一、标题
% ─────────────────────────────────────────────────────────────────────

% 参考 EMSWOA 标题模式：
%   [Mechanism 1] + [Mechanism 2] + for [application domain]

% KDS-EL 标题（最终版）：
\title{Knowledge-guided Dual-Space Evolutionary Learning for
High-Dimensional Bioinformatics Feature Selection}

% 备用（更强调双目标特性）：
% \title{A Bi-objective-regulated Dual-Space Evolutionary Framework
% with Knowledge-guided Elite Refinement for High-Dimensional Feature Selection}


% ─────────────────────────────────────────────────────────────────────
% 二、Highlights（5条，动词开头，单行内）
% ─────────────────────────────────────────────────────────────────────

\begin{highlights}
\item Propose a bi-objective-regulated dual-space evolutionary learning
      framework (KDS-EL) for high-dimensional bioinformatics feature selection.
\item Design a Diversity-regulated Selection (DRS) operator that elevates
      NHD into a first-class Pareto selection objective to prevent structural collapse.
\item Develop a Knowledge-guided Elite Refinement (KER) mechanism that
      accumulates improvement evidence online to direct targeted bit-flip perturbations.
\item Demonstrate synergistic gain of DRS and KER on 8 of 9 benchmark datasets.
\item Achieve superior classification accuracy and maximum feature reduction
      against eight state-of-the-art competitors including PSO-CSM (2025) and IGPSO (2024).
\end{highlights}


% ─────────────────────────────────────────────────────────────────────
% 三、Keywords（5个）
% ─────────────────────────────────────────────────────────────────────

Feature selection; Evolutionary computation; Population diversity;
Knowledge-guided learning; Bioinformatics


% ─────────────────────────────────────────────────────────────────────
% 四、Abstract（~200词，严格参照 EMSWOA 结构）
% ─────────────────────────────────────────────────────────────────────

% 句式结构（对照 EMSWOA）：
%   Sentence 1-2: 问题重要性 + 核心困难
%   Sentence 3:   提出 KDS-EL
%   Sentence 4-5: DRS 贡献
%   Sentence 6-7: KER 贡献
%   Sentence 8:   实验规模
%   Sentence 9:   核心结果数字

\begin{abstract}
Feature selection is essential for reducing dimensionality and improving
classifier performance in high-dimensional bioinformatics data, where
the exponential search space makes evolutionary methods highly
susceptible to structural population collapse and knowledge waste.
To address these coupled challenges, we propose \textbf{KDS-EL}
(\textit{Knowledge-guided Dual-Space Evolutionary Learning}), a novel
bi-objective-regulated evolutionary framework that simultaneously
optimises a dual-population PSO backbone across two complementary
search spaces.
First, we design a \textbf{Diversity-regulated Selection} (DRS) operator
that elevates Normalised Hamming Distance (NHD) from a post-hoc
diagnostic metric to a first-class Pareto selection objective,
providing a provable structural diversity lower bound and preventing
premature population convergence.
Second, we develop a \textbf{Knowledge-guided Elite Refinement} (KER)
mechanism that accumulates improvement-event evidence from the
evolutionary trajectory into a feature-importance weight vector,
and uses this online-learned knowledge to direct targeted bit-flip
perturbations of elite solutions without additional fitness evaluations.
The two mechanisms are synergistically coupled: DRS's structural
diversity ensures heterogeneous improvement signals for KER, while KER's
directed perturbations enrich the DRS selection pool.
Extensive experiments on nine high-dimensional bioinformatics datasets
($d \in [927,\,26{,}877]$) against eight state-of-the-art competitors
demonstrate that KDS-EL achieves the highest average Friedman rank
of 1.44 (Friedman $\chi^2=52.03$, $p=1.66\times10^{-8}$),
reduces selected feature counts by 75\%--88\% relative to the base
framework, and outperforms recent methods including PSO-CSM (2025)
and IGPSO (2024) across all nine datasets with statistical significance.
Ablation analysis confirms positive synergistic gain on 8 of 9 datasets,
validating the Dual-Space design rationale.
\end{abstract}


% ─────────────────────────────────────────────────────────────────────
% 五、完整论文结构（节-小节-段落级）
% ─────────────────────────────────────────────────────────────────────

\section{Introduction}
% 参照 EMSWOA §1 结构，约 800-1000 词，5 段
%
% §1 P1: 应用背景 — 高维组学数据的重要性
%   "With the rapid advances in high-throughput sequencing, gene expression datasets
%    routinely comprise tens of thousands of features over fewer than 200 samples..."
%
% §1 P2: Limitation 1 — 结构坍缩 (对应 EMSWOA 的"单种群"局限)
%   "Most EC-based FS methods treat the search as a single-objective scalar problem,
%    leaving the structural geometry of the binary solution space unexploited..."
%
% §1 P3: Limitation 2 — 知识浪费 (对应 EMSWOA 的"无效概率翻转"局限)
%   "Furthermore, the cumulative evidence of which feature transitions proved
%    beneficial is discarded at every generation..."
%
% §1 P4: 提出 KDS-EL + 双空间定义
%   "To address these limitations, we propose KDS-EL..."
%   定义 Fitness Space + Structural Space
%   说明 DRS + KER 的作用
%   强调协同效应
%
% §1 P5: Contributions (用 enumerate，参照 EMSWOA 5 bullet highlights)
%   (i)  提出 KDS-EL 框架
%   (ii) DRS 算子 + 可证明多样性下界
%   (iii)KER 机制 + 在线知识积累
%   (iv) 综合实验验证
%
% §1 末: 论文结构说明
%   "The remainder of this paper is organised as follows. Section 2 reviews..."


\section{Related Work}

\subsection{Feature Selection Methods}
% 综述 filter / wrapper / embedded 三类
% 重点说明 wrapper + EC 的优势和现有局限
% 引用：~8-10 篇，涵盖 BPSO, FWPSO, SaWDE, MEL, SAR-MTPSO
% 结尾句：引出多样性问题

\subsection{Population Diversity in Evolutionary Feature Selection}
% 说明多样性维护的重要性
% 现有方法：多子群、种群重组、拥挤距离
% NHD 作为诊断指标 vs 本文将其升格为选择目标
% 关键句："To the best of our knowledge, KDS-EL is the first to
%           incorporate NHD as a Pareto selection objective..."

\subsection{Knowledge-guided and Multi-task Evolutionary FS}
% 综述 IGPSO (2024) — MLP预训练
% 综述 SAR-MTPSO (2025) — 多任务异构目标
% 综述 PSO-CSM (2025) — 综合评分
% 引出 KER 的不同之处：无需预训练、在线累积、单一可解释参数
% 处理 SAR-MTPSO：双任务异构目标框架与 KDS-EL 单目标框架不可直接比较

% 注意：参照 EMSWOA Table 1 格式，加一个文献分析表
% Table 1: Analysis of recent PSO-based FS algorithms
%   列：Algorithm, Year, High-dim, Diversity mechanism, Knowledge-guided
%   行：BPSO, FWPSO, MEL, IGPSO, PSO-CSM, KDS-EL


\section{The Proposed KDS-EL Framework}
% 参照 EMSWOA §3 结构

\subsection{Problem Formulation}
% 标准 FS 问题公式
% fitness = alpha * ErrorRate + (1-alpha) * (#Selected/#All)
% alpha = 0.9 (与 EMSWOA 的 alpha 定义一致，但我们用 [0.9, 0.1])

\subsection{Dual-Population PSO Backbone}
% 说明 Sub1 (PSO探索) + Sub2 (知识引导采样)
% 引用 MEL [CITE_MEL]
% 强调这是已有基础，DRS 和 KER 是本文贡献

\subsection{Diversity-regulated Selection (DRS)}
% 对应 EMSWOA §3.2 的详细程度
%
% Step 1: NHD 定义和池化
%   公式：NHD(x_i, x_j) = AHD / max(|x_i|+|x_j|-2·inter, 1)
%   公式：pool = Pbest ∪ P_new (size = 2N)
%
% Step 2: 双目标 Pareto 排序
%   目标1: min fitness (classification error + feature ratio)
%   目标2: max structural diversity (NHD)
%   使用 NSGA-II 非支配排序
%
% Step 3: 拥挤距离 + 自适应权重
%   alpha_cd(t) = (t/T)^2  — 早期多样性优先，后期精度优先
%   score = alpha_cd * CD_norm + (1 - alpha_cd) * NHD_norm
%
% Proposition 1: 多样性下界
%   "For t < T, the DRS operator guarantees NHD > 0 in the selected population."
%   Proof sketch: alpha_cd < 1 → NHD term always positive in score
%
% Algorithm 1: DRS 伪代码（参照 EMSWOA Algorithm 1 格式）

\subsection{Knowledge-guided Elite Refinement (KER)}
% 对应 EMSWOA §3.3 的详细程度
%
% Step 1: 权重更新（BEFORE Pbest 更新）
%   improvement event: delta = fitP[i] - fit_new[i]
%   weight[feature added] += delta
%   weight[feature removed] -= delta
%
% Step 2: 精英扰动
%   prob_random = rand_flip_p = 0.10 (ergodicity保证)
%   prob_guided: scores = weight * (1-elite) - weight * elite + noise
%   mutant = flip top-k bits
%
% Proposition 2: 遍历性保证
%   "With probability p_r > 0 at each iteration, KER explores any binary mask."
%
% Algorithm 2: KER 伪代码（参照 EMSWOA Algorithm 2 格式）

\subsection{Complete KDS-EL Algorithm}
% Algorithm 3: 完整伪代码
% 包含 DRS + KER 的完整流程
% 参数表格（对应 EMSWOA Table 3 格式）

\subsection{Fitness Function}
% 标准公式（与 EMSWOA §3.4 相同格式）
% 说明 K-fold CV 设置：K=5 (EMSWOA用K=10，说明我们的选择理由)

\subsection{Time Complexity Analysis}
% 参照 EMSWOA §3.5 格式
% 逐组件分析
% 结论：O(T·(2N+k)·C_f + T·N²·d)
% 说明实际开销与 BPSO/SaWDE 相当


\section{Experimental Results and Discussion}

\subsection{Experimental Setup}

\subsubsection{Datasets}
% Table 2: 9 datasets (Dataset, #Features, #Instances, #Classes, Source)
% 参照 EMSWOA Table 2 格式，但我们只有9个 → 补充说明为何选这9个
% 可考虑加 3-4 个 UCI 数据集扩充到 12-13 个（提高说服力）

\subsubsection{Comparison Algorithms}
% 8 algorithms: BPSO, SaWDE, FWPSO, MEL, KDS-EL/DRS, KDS-EL/KER, PSO-CSM, IGPSO
% Table 3: Parameter settings（参照 EMSWOA Table 3）

\subsubsection{Evaluation Protocol}
% 10 independent runs (EMSWOA没提runs数，我们明确说明)
% StratifiedKFold(5, shuffle=True, random_state=42)
% Seeds: 42 + run × 37
% Metrics: Accuracy (mean ± std), Num_Features, Avg_NHD, Time_sec

\subsection{Experimental Analysis}

\subsubsection{Classification Accuracy Comparison}
% Table 4: 9方法×9数据集，mean ± std，bold best
% 参照 EMSWOA Table 4 格式（bold最优，underline次优）
% Fig 1: 精度收敛曲线（参照 EMSWOA Fig 7 格式）
%   → 分9个子图，每图9条线
%   ⚠️ 注意：EMSWOA 画的是 accuracy 收敛曲线，我们目前只有 fitness 曲线
%      需要在 npy 中同时保存 accuracy 曲线（或转换）

\subsubsection{Feature Subset Size Comparison}
% Table 5: 特征数对比（参照 EMSWOA Table 5）
% Fig 2: 特征数柱状图（选4-5个代表性数据集对比）

\subsubsection{Fitness Value Analysis}
% Table 6: fitness 值 (mean ± std)（参照 EMSWOA Table 6）
% Fig 3: Fitness 收敛曲线（参照 EMSWOA Fig 8）

\subsubsection{Statistical Significance Tests}
% Table 7: Wilcoxon p 值表（参照 EMSWOA Table 7）
%   格式：行=数据集，列=对比方法，值=p值
%   粗体：p > 0.05（即不显著），下划线：0.01 < p < 0.05
% + Friedman 检验结果（EMSWOA 没有 → 这是 KDS-EL 的加分项）
% + Holm-Bonferroni 校正（可选，放在 Appendix）

\subsubsection{Ablation Study}
% Table 8: 4个变体对比（参照 EMSWOA Table 8 格式）
%   行=数据集，列=MEL / KDS-EL/DRS / KDS-EL/KER / KDS-EL
%   给出 ACC ± STD 和 Num_Features
% Fig 4: 消融实验双柱图（已有 fig3_ablation.png）
% 协同增益表：Δ_full - Δ_DRS - Δ_KER（新增，EMSWOA 没有 → 加分项）

\subsubsection{Stability Analysis}
% Fig 5: 箱线图（已有 fig4_boxplot.png）
% 简短讨论：std 指标，DRS 维持多样性 → 稳定性更高

\subsubsection{Convergence Analysis}
% Fig 6: 收敛曲线（已有 fig10a_convergence）
% Fig 7: NHD 多样性曲线（已有 fig10b_nhd_curves）

\subsubsection{Computational Cost Analysis}
% Fig 8: Runtime 对比（已有 fig8_runtime.png 双面板版）
% Fig 9: 气泡图 accuracy-feature-runtime（已有 fig_bubble_tradeoff）
% Table 9: 时间复杂度表（已有 complexity_section.tex）

\subsubsection{Discussion: ALL3 Exception}
% 约 400 词
% n=36, d=7129 极端小样本 → KNN 5折每折7样本 → 信号噪声比低
% KER 的 improvement feedback 不可靠 → 权重噪声累积
% 结论：KDS-EL 在 n < 40 时性能下降，这是方法的边界条件

\subsection{Discussion}
% 参照 EMSWOA §4.2.8
% 综合分析两个机制的作用
% 强调协同效应
% 说明应用价值（biomarker discovery）
% 承认局限性（小样本，ALL3）


\section{Conclusions and Future Work}
% 参照 EMSWOA §5 结构，约 300-400 词
%
% P1: 问题 + 提出 KDS-EL（一句重述）
% P2: DRS 的贡献 + 实验证明
% P3: KER 的贡献 + 实验证明
% P4: 整体结果（rank 1.44, 8/9 协同, 75-88% 特征减少）
% P5: 局限性（极端小样本 n<40，ALL3 例外）
% P6: Future work（3条）
%   - 扩展到多标签特征选择
%   - 动态 d 自适应（可变长度特征）
%   - 结合深度学习知识转移（参照 EMSWOA 结尾）


% ─────────────────────────────────────────────────────────────────────
% 六、图表清单（最终版）
% ─────────────────────────────────────────────────────────────────────

% 参照 EMSWOA 的图表数量（10 figures, 10 tables）

% FIGURES（共 10 张）：
% Fig 1:  KDS-EL 框架流程图 + 特征选择示例（参照 EMSWOA Fig 1+2 合并）
%         → 需要新画：包含 DRS + KER 的完整流程图
% Fig 2:  DRS 机制示意图（Pareto 前沿 + NHD 结构）
%         → 需要新画
% Fig 3:  KER 机制示意图（权重更新 + 精英扰动流程）
%         → 需要新画
% Fig 4:  准确率对比（3×3，9方法）→ 已有 fig1_accuracy.png
% Fig 5:  特征数对比（3×3，log）→ 已有 fig2_features.png
% Fig 6:  消融实验（双面板）→ 已有 fig3_ablation.png
% Fig 7:  收敛曲线（准确率，2×3，6代表数据集）→ 需修改（目前是fitness）
% Fig 8:  NHD 多样性曲线 → 已有 fig10b_nhd_curves
% Fig 9:  气泡图（accuracy-feature-runtime）→ 已有 fig_bubble_tradeoff
% Fig 10: Wilcoxon 显著性热图 → 已有 fig9_wilcoxon.png

% TABLES（共 8 张）：
% Table 1: 近期 PSO-based FS 文献分析表（参照 EMSWOA Table 1）
%           → 需要新写：列=Algorithm/Year/High-dim/Diversity/Knowledge
% Table 2: 数据集详情（Dataset, Features, Instances, Classes）
%           → 已知数据
% Table 3: 参数设置（所有算法）
%           → 已知参数
% Table 4: 准确率对比（mean ± std，bold最优）
%           → 需要生成 LaTeX 表格
% Table 5: 特征数对比（bold最优）
%           → 需要生成 LaTeX 表格
% Table 6: Wilcoxon p 值表（参照 EMSWOA Table 7）
%           → 需要生成
% Table 7: 消融实验结果（4变体，ACC+Features）
%           → 需要生成
% Table 8: 时间复杂度对比表
%           → 已有 complexity_section.tex


% ─────────────────────────────────────────────────────────────────────
% 七、与 EMSWOA 的差距分析和补充建议
% ─────────────────────────────────────────────────────────────────────

% EMSWOA 有 KDS-EL 没有的：
%
% ① NON-EC 对比（Filter 方法）
%    EMSWOA 对比了 S2DFS, EGCFS, MSVM-RFE 等非进化方法
%    建议：加入 ReliefF, LASSO, SVM-RFE 对比
%    工作量：只需从文献引用已有数字，不需要重跑
%    价值：大幅提升论文完整性，审稿人会问为什么不对比 filter 方法
%
% ② 更多数据集（19 vs 9）
%    建议：从 EMSWOA 的19个数据集里选4-5个 UCI 通用数据集补充
%    如：Yale(1024d), Brain_Tumor(5920d), SMK-CAN-187(19993d)
%    工作量：在 Kaggle 上重跑4-5个数据集，每个约2-3小时
%    价值：从9→13个数据集，说服力大幅提升
%
% ③ 准确率收敛曲线（不是fitness曲线）
%    EMSWOA 展示的是 accuracy 随迭代变化（不是 fitness）
%    更直观，审稿人更容易理解
%    建议：在 npy 中同时记录每次迭代的 Gbest 对应的 accuracy
%
% KDS-EL 有 EMSWOA 没有的（优势）：
%
% ① Friedman 检验（EMSWOA 只有 Wilcoxon）
% ② 协同效应定量分析（Δ_full > Δ_DRS + Δ_KER）
% ③ NHD 多样性曲线（直接验证 DRS 效果）
% ④ 气泡图（三维 trade-off 可视化）
% ⑤ 双空间理论框架（更系统的理论定位）


% ─────────────────────────────────────────────────────────────────────
% 八、写作时间规划
% ─────────────────────────────────────────────────────────────────────

% 第一周：
%   Day 1-2: 画 Fig 1 (框架图) + Fig 2 (DRS示意图) + Fig 3 (KER示意图)
%   Day 3:   Table 1 (文献分析表) + Table 2-3 (数据集+参数)
%   Day 4-5: 生成所有 LaTeX 表格 (Table 4-8)
%   Day 6-7: 写 Methodology §3 完整版（已有框架，填充细节）

% 第二周：
%   Day 1-2: 写 Related Work §2
%   Day 3-4: 写 Introduction §1 完整版
%   Day 5-6: 写 Experiments §4 (有数字有图，主要是语言组织)
%   Day 7:   写 Conclusion §5

% 第三周：
%   Day 1-2: 决定是否补充 NON-EC 对比（引用文献数字）
%   Day 3-4: 全文通读修改
%   Day 5:   格式检查（ESWA/ASC 模板套用）
%   Day 6-7: 最终审阅 + 投稿
